% Bài: L12C1 Bài 2. Nội năng. Định luật I của nhiệt động lực học · Nội năng và truyền nhiệt
% ID: L12C1B3-121-TN
% ID: L12C1B3-121-TN
% Mức: NB
% ID: L12C1B3-121-TN


% Mức: TH
% Mức: TH
% ID: L12C1B3-123-TN
% ID: L12C1B3-123-TN
% Mức: TH
% ID: L12C1B3-123-TN
% ID: L12C1B3-123-TN
% Mức: TH
% ID: L12C1B2-02-TN
%=== Câu 3 ===%
\dangbt{Nội năng và truyền nhiệt}
\begin{ex}
% ID: L12C1B2-02-TN
% Mức: TH
	Cho hai vật $A$ và $B$ tiếp xúc nhau. Nhiệt chỉ tự truyền từ vật $A$ sang vật $B$ khi
	\choice 
	{$A$ và $B$ là hai vật rắn}
	{nhiệt độ của $A$ và của $B$ bằng nhau}
	{\True nhiệt độ của $A$ lớn hơn nhiệt độ của $B$}
	{khối lượng của $A$ lớn hơn khối lượng của $B$}
	\loigiai{
		Theo chiều truyền năng lượng nhiệt: Năng lượng nhiệt sẽ tự truyền từ vật có nhiệt độ cao hơn sang vật có nhiệt độ thấp hơn,. Quá trình truyền nhiệt kết thúc khi hai vật có cùng nhiệt độ (trạng thái cân bằng nhiệt).
		
		Vì vậy, nhiệt chỉ tự truyền từ vật $A$ sang vật $B$ khi nhiệt độ của $A$ lớn hơn nhiệt độ của $B$.
	}
\end{ex}

% Mức: TH
% Mức: TH
% ID: L12C1B3-126-TN
% ID: L12C1B3-126-TN
% Mức: TH
% ID: L12C1B3-126-TN
% ID: L12C1B3-126-TN
% Mức: TH
% ID: L12C1B2-03-TN
%=== Câu 6 ===%
\begin{ex}
% ID: L12C1B2-03-TN
% Mức: TH
	 Nếu làm tăng nhiệt độ của một hệ mà không làm thay đổi thể tích của nó thì nội năng của hệ sẽ
	 \choice 
	 {\True tăng lên}
	 {không thay đổi}
	 {giảm xuống}
	 {tăng lên sau đó giảm xuống}
	\loigiai{
	Khi nhiệt độ tăng thì động năng
	}
\end{ex}

% Mức: NB
% Mức: NB
% ID: L12C1B3-127-TN
% ID: L12C1B3-127-TN
% Mức: NB
% ID: L12C1B3-127-TN
% ID: L12C1B3-127-TN
% Mức: NB
% ID: L12C1B2-04-TN
%=== Câu 7 ===%
\begin{ex}
% ID: L12C1B2-04-TN
% Mức: NB
	Phần năng lượng nhiệt mà vật này truyền cho vật kia hoặc vật này nhận từ vật kia gọi là
	\choice 
	{nội năng}
	{thế năng}
	{nhiệt độ}
	{\True nhiệt lượng}
	\loigiai{
		Phần năng lượng nhiệt truyền từ vật có nhiệt độ cao hơn sang vật có nhiệt độ thấp hơn (hay phần năng lượng nhiệt trao đổi giữa hai vật) trong quá trình truyền nhiệt được gọi là nhiệt lượng,.
	}
\end{ex}

% Mức: TH
% Mức: TH
% ID: L12C1B3-133-DS
% ID: L12C1B3-133-DS
% Mức: TH
% ID: L12C1B3-133-DS
% ID: L12C1B3-133-DS
% Mức: TH
% ID: L12C1B2-05-DS
%=== Câu 13 ===%
\begin{ex}
% ID: L12C1B2-05-DS
% Mức: TH
	Có hai cốc nước $A$ và $B$ chứa cùng một lượng nước ở nhiệt độ phòng. Người ta thả một viên nước đá vào cốc $A$ và nhúng cốc $B$ vào một bình chứa nước nóng.
	\choiceTF
	{\True Cốc $B$ nhận nhiệt lượng từ nước ở bình và nhiệt độ nước trong cốc tăng lên}
	{Nhiệt độ nước ở cốc $A$ giảm vì nhận nhiệt lượng từ nước đá}
	{\True Nhiệt lượng không thể tự truyền từ nước ở cốc $B$ vào nước ở bình chứa nó}
	{\True Khi nhiệt độ nước ở cốc $B$ là $t_2 =62^\circ C$ (theo thang nhiệt Celsius) thì theo thang nhiệt độ Kelvin là $T_2 = 335\ K$}
	\loigiai{
		\begin{itemchoice}
			\itemch Phát biểu này đúng. Cốc $B$ được đặt trong nước nóng, nên sẽ nhận nhiệt từ nước nóng làm nhiệt độ nước trong cốc tăng lên.
			\itemch Phát biểu này sai. Nhiệt độ nước ở cốc $A$ giảm vì \textbf{truyền} nhiệt lượng cho nước đá (nước đá nhận nhiệt lượng từ nước).
			\itemch  Phát biểu này đúng. Nhiệt lượng không thể tự truyền từ vật lạnh sang vật nóng mà không có tác động từ bên ngoài.
			\itemch  Phát biểu này đúng. Công thức chuyển đổi từ độ Celsius sang Kelvin là $T(K) = t(^\circ C) + 273$. Do đó, $T_2 = 62 + 273 = 335\ K$.
		\end{itemchoice}
	}
\end{ex}
